% Part: first-order-logic
% Chapter: beyond
% Section: second-order-logic

\documentclass[../../../include/open-logic-section]{subfiles}

\begin{document}

\olfileid[hi]{fol}{byd}{sol}

\olsection{द्वितीय-कोटि तर्कशास्त्र}

द्वितीय-कोटि तर्कशास्त्र की भाषा केवल व्यष्टियों के किसी !!a{domain} पर ही
नहीं, बल्कि उस !!{domain} पर संबंधों के ऊपर भी परिमाणीकरण करने देती है। कोई
प्रथम-कोटि भाषा $\Lang{L}$ दी हो, तो प्रत्येक $k$ के लिए !!{variable}s का एक
भंडार जोड़ते हैं; इन चरों को $R$ लिखते हैं और वे $k$-पदी संबंधों पर विचरते हैं। उन चरों पर परिमाणीकरण की अनुमति
दी जाती है। यदि $R$ किसी $k$-पदी संबंध का !!a{variable} और $t_1$,
\dots,~$t_k$ सामान्य (प्रथम-कोटि) पद हैं, तो $\Atom{R}{t_1,\dots,t_k}$ एक
परमाण्विक !!{formula} है। अन्यथा !!{formula}s का समुच्चय प्रथम-कोटि
तर्कशास्त्र की तरह ही, द्वितीय-कोटि परिमाणीकरण के अतिरिक्त उपवाक्यों सहित,
परिभाषित होता है। ध्यान दें कि हमारे पास केवल प्रथम-कोटि पदों के लिए
!!{identity} है: यदि $R$ और $S$ संबंध-!!{variable}s के ऐसे उदाहरण हैं जिनकी
समान पदिता~$k$ है, तो $\eq[R][S]$ को निम्न का संक्षेप परिभाषित कर सकते हैं:
\[
\lforall[x_1 \dots][\lforall[x_k][(\Atom{R}{x_1, \dots, x_k} \liff
  \Atom{S}{x_1, \dots, x_k})]].
\]

द्वितीय-कोटि तर्कशास्त्र के नियम परिमाणक-नियमों को नए द्वितीय-कोटि चरों तक
विस्तृत करते हैं। किंतु यहाँ सावधानी से बताना पड़ता है कि ये चर~$\Lang{L}$
के !!{predicate}s, और अधिक सामान्य रूप से~$\Lang{L}$ के !!{formula}s, के
साथ कैसे अंतःक्रिया करते हैं। कम-से-कम संबंध-चर पद माने जाते हैं, अतः हमारे
पास इस रूप के अनुमान होते हैं:
\[
!A(R) \Proves \lexists[R][!A(R)]
\]
पर यदि $\Lang{L}$ अंकगणित की भाषा है जिसमें अचर संबंध-चिह्न~$<$ है, तो हम
निम्न अनुमान को भी वैध अपेक्षित करेंगे:
\[
x < y \Proves \lexists[R][\Atom{R}{x,y}]
\]
अथवा किसी दिए !!{formula}~$!A$ के लिए,
\[
!A(x_1, \dots, x_k) \Proves \lexists[R][\Atom{R}{x_1,\dots,x_k}]
\]
अधिक सामान्यतः हम निम्न रूप के अनुमानों की अनुमति देना चाह सकते हैं:
\[
\Subst{!A}{\lambd[\vec x][!B(\vec x)]}{R} \Proves
\lexists[R][!A]
\]
जहाँ $\Subst{!A}{\lambd[\vec x][!B(\vec x)]}{R}$,
$\Obj{R}{t_1,\dots,t_k}$ रूप के प्रत्येक परमाण्विक !!{formula} को~$!A$ में
$!B(t_1, \dots, t_k)$ से बदलने पर मिला परिणाम निरूपित करता है। यह अंतिम
नियम एक {\em परिग्रहण-योजना} रखने के समतुल्य है, अर्थात् इस रूप का अभिगृहीत:
\[
\lexists[R][\lforall[x_1, \dots, x_k][(!A(x_1, \dots, x_k) \liff
\Atom{R}{x_1, \dots, x_k})]],
\]
द्वितीय-कोटि भाषा के प्रत्येक ऐसे !!{formula}~$!A$ के लिए एक, जिसमें $R$
मुक्त चर नहीं है। (अभ्यास: दिखाइए कि यदि $R$ को~$!A$ में आने दिया जाए, तो
यह योजना असंगत है!)

तर्कशास्त्री जब ``द्वितीय-कोटि तर्कशास्त्र के अभिगृहीत'' कहते हैं, तो उनका
आशय सामान्यतः परिग्रहण-योजना सहित द्वितीय-कोटि परिमाणक-नियमों द्वारा
प्रथम-कोटि तर्कशास्त्र के न्यूनतम विस्तार से होता है। किंतु इन अभिगृहीतों
और नियमों के दुर्बलतर उपतंत्रों का अध्ययन भी प्रायः रोचक होता है। उदाहरणतः
ध्यान दें कि अपनी पूर्ण व्यापकता में परिग्रहण की अभिगृहीत-योजना
\emph{अविधेयात्मक} है: वह ऐसे संबंध $\Atom{R}{x_1, \dots, x_k}$ के अस्तित्व
का कथन करने देती है जो द्वितीय-कोटि परिमाणकों वाले किसी !!a{formula} से
``परिभाषित'' है; और ये परिमाणक ऐसे सभी संबंधों के समुच्चय पर विचरते हैं---
ऐसा समुच्चय जिसमें स्वयं $R$ भी है!{} बीसवीं शताब्दी के आरंभ के आसपास रसेल
के विरोधाभास पर एक सामान्य प्रतिक्रिया ऐसी परिभाषाओं को दोषी मानना और
गणित की नींव विकसित करते समय उनसे बचना था। यदि !!{formula}~$!A$ में
द्वितीय-कोटि परिमाणकों का प्रयोग निषिद्ध कर दें, तो परिग्रहण का
{\em विधेयात्मक} रूप मिलता है, जो कुछ दुर्बल है।

अर्थगत दृष्टि से द्वितीय-कोटि !!{structure} को भाषा की एक प्रथम-कोटि
!!{structure} और उस !!{domain} पर संबंधों के ऐसे समुच्चय से बना मान सकते हैं
जिस पर द्वितीय-कोटि परिमाणक विचरते हैं (अधिक सटीकतः प्रत्येक $k$ के लिए
पदिता~$k$ वाले संबंधों का एक समुच्चय होता है)। निस्संदेह, यदि
!!{derivation}-तंत्र में परिग्रहण सम्मिलित है, तो एक अतिरिक्त आवश्यकता है
कि ``द्वितीय-कोटि भाग'' में परिग्रहण-अभिगृहीतों को संतुष्ट करने के लिए पर्याप्त
संबंध हों---अन्यथा !!{derivation}-तंत्र सुदृढ़ नहीं है!{} पर्याप्त संबंध
सुनिश्चित करने की एक सरल रीति यह है कि द्वितीय-कोटि भाग में प्रथम-कोटि भाग
पर \emph{सभी} संबंध ले लें। ऐसी !!a{structure} को \emph{पूर्ण} कहा जाता है
और एक अर्थ में वही भाषा की वास्तविक ``अभीष्ट !!{structure}'' है। यदि हमारे
विचार में आने वाली !!{structure}s केवल पूर्ण हों, तो \emph{पूर्ण} द्वितीय-कोटि
अर्थविज्ञान मिलता है। उस स्थिति में किसी !!{structure} का विनिर्देशन केवल
प्रथम-कोटि भाग के विनिर्देशन में सिमट जाता है, क्योंकि द्वितीय-कोटि भाग की
विषयवस्तु उससे अंतर्निहित रूप से मिलती है।

सारतः द्वितीय-कोटि तर्कशास्त्र की बात में कुछ संदिग्धता है।
!!{derivation}-तंत्र की दृष्टि से हमारा आशय इनमें से कोई हो सकता है:
\begin{enumerate}
\item कुछ परिग्रहण-अभिगृहीतों सहित एक ``न्यूनतम'' द्वितीय-कोटि
  !!{derivation}-तंत्र।
\item पूर्ण परिग्रहण वाला ``मानक'' द्वितीय-कोटि !!{derivation}-तंत्र।
\end{enumerate}
अर्थविज्ञान की दृष्टि से इनमें से किसी में रुचि हो सकती है:
\begin{enumerate}
\item ``दुर्बल'' अर्थविज्ञान, जिसमें !!a{structure} एक प्रथम-कोटि भाग और
  परिग्रहण-अभिगृहीतों को संतुष्ट करने जितने बड़े द्वितीय-कोटि भाग से बनती है।
\item ``मानक'' द्वितीय-कोटि अर्थविज्ञान, जिसमें विचाराधीन !!{structure}s
  केवल पूर्ण होती हैं।
\end{enumerate}
जब तर्कशास्त्री अपने अभीष्ट !!{derivation}-तंत्र या अर्थविज्ञान को स्पष्ट
नहीं करते, तो सामान्यतः प्रत्येक सूची की दूसरी मद का निर्देश करते हैं। इस
अर्थविज्ञान का लाभ यह है कि, जैसा हम देखेंगे, इससे अनेक स्वाभाविक गणितीय
संरचनाओं के श्रेणीपरक वर्णन मिलते हैं; साथ ही !!{derivation}-तंत्र काफी
प्रबल और इस अर्थविज्ञान के लिए सुदृढ़ है। दोष यह है कि !!{derivation}-तंत्र
अर्थविज्ञान के लिए \emph{पूर्ण नहीं} है; वस्तुतः प्रभावी ढंग से दिया गया
\emph{कोई भी} !!{derivation}-तंत्र पूर्ण द्वितीय-कोटि अर्थविज्ञान के लिए
पूर्ण नहीं है। दूसरी ओर हम देखेंगे कि !!{derivation}-तंत्र दुर्बलतर
अर्थविज्ञान के लिए \emph{पूर्ण है}; इससे निष्पन्न होता है कि यदि कोई वाक्य
प्रमाण्य नहीं है, तो \emph{कोई} !!{structure}---जो अनिवार्यतः पूर्ण न
हो---ऐसी अवश्य है जिसमें वह असत्य है।

द्वितीय-कोटि तर्कशास्त्र की भाषा बहुत अभिव्यंजक है। एकपदी संबंधों की पहचान
!!{domain} के उपसमुच्चयों से की जा सकती है और इस कारण इन समुच्चयों पर
परिमाणीकरण भी किया जा सकता है; जैसे, प्राकृतिक संख्याओं के लिए आगमन को एक
ही अभिगृहीत में व्यक्त कर सकते हैं:
\[
\lforall[R][((\Atom{R}{\Obj{0}} \land \lforall[x][(\Atom{R}{x} \lif
    \Atom{R}{x'})]) \lif \lforall[x][\Atom{R}{x}])].
\]
यदि अंकगणित की भाषा में $\Obj 0, \Obj \prime, +,
\times$ और $<$ चिह्न लें,
तो उनके व्यवहार के वर्णन के लिए निम्न अभिगृहीत जोड़ सकते हैं:
\begin{enumerate}
\item $\lforall[x][\lnot x' = \Obj 0]$
\item $\lforall[x][\lforall[y][(s(x) = s(y) \lif x = y)]]$
\item $\lforall[x][(x + \Obj 0) = x]$
\item $\lforall[x][\lforall[y][(x + y') = (x + y)']]$
\item $\lforall[x][(x \times \Obj 0) = \Obj 0]$
\item $\lforall[x][\lforall[y][(x \times y') = ((x \times y) + x)]]$
\item $\lforall[x][\lforall[y][(x < y \liff \lexists[z][y = (x + z')])]]$
\end{enumerate}
यह दिखाना कठिन नहीं है कि यदि हम पूर्ण द्वितीय-कोटि अर्थविज्ञान का प्रयोग
कर रहे हैं, तो ये अभिगृहीत ऊपर के आगमन-अभिगृहीत के साथ अंकगणित के मानक
प्रतिरूप !!{structure}~$\Struct{N}$ का श्रेणीपरक वर्णन देते हैं। ऐसी कोई
!!{structure}~$\Struct{M}$ दी हो जिसमें ये अभिगृहीत सत्य हैं, तो फलन~$f$ को
$\Nat$ से $\Struct{M}$ के !!{domain} में, $\Nat$ पर सामान्य पुनरावर्तन
द्वारा, ऐसे परिभाषित करें कि $f(0) = \Assign{\Obj 0}{M}$ और
$f(x+1) = \Assign{\prime}{M}(f(x))$। $\Nat$ पर सामान्य आगमन और~$\Struct M$ में
अभिगृहीत (1) तथा~(2) सत्य होने के तथ्य से देखते हैं कि $f$ !!{injective}
है। $f$ को !!{surjective} देखने के लिए~$P$ को~$\Domain{M}$ के उन अवयवों का
समुच्चय लें जो~$f$ के परिसर में हैं। क्योंकि $\Struct M$ पूर्ण है, $P$
द्वितीय-कोटि !!{domain} में है।~$f$ की रचना से जानते हैं कि
$\Assign{\Obj
  0}{M}$,~$P$ में है और $P$,~$\Assign{\prime}{M}$ के अधीन संवृत
है। $\Struct M$ में (विशेषतः~$P$ के लिए) आगमन-अभिगृहीत सत्य होने का तथ्य
प्रत्याभूत करता है कि $P$,~$\Struct M$ के पूरे प्रथम-कोटि !!{domain} के
बराबर है। इससे सिद्ध होता है कि $f$ !!a{bijection} है। यह दिखाना भी अधिक
कठिन नहीं कि $f$ एक समाकृतिक प्रतिचित्रण है; इसके लिए $\Nat$ पर सामान्य आगमन का बार-बार
प्रयोग करें।

समुच्चय-सैद्धांतिक पदों में कोई फलन केवल एक विशेष प्रकार का संबंध है; जैसे,
एकपदी फलन~$f$ की पहचान ऐसे द्विपदी संबंध~$R$ से की जा सकती है जो
$\lforall[x][\lexists![y][R(x,y)]]$ को संतुष्ट करता है। फलतः फलनों पर भी परिमाणीकरण कर सकते
हैं। पूर्ण अर्थविज्ञान का प्रयोग करके ऐसा वर्ग परिभाषित कर सकते हैं जिसकी
सदस्य !!{structure}s अनंत हों: ये वे !!{structure}s $\Struct M$ हैं जिनके लिए $\Struct M$
के !!{domain} से उसके ही किसी उचित उपसमुच्चय में एक एकैकी फलन है:
\[
\lexists[f][(\lforall[x][\lforall[y][(\eq[f(x)][f(y)] \lif \eq[x][y])]]
  \land \lexists[y][\lforall[x][\eq/[f(x)][y]]])].
\]
तब इस वाक्य का निषेध उस वर्ग को परिभाषित करता है जिसकी सदस्य
!!{structure}s परिमित हैं।

इसके अतिरिक्त किसी रैखिक क्रम की परिभाषा में निम्न जोड़कर सुक्रमों का वर्ग
परिभाषित किया जा सकता है:
\[
\lforall[P][(\lexists[x][\Atom{P}{x}] \lif \lexists[x][(\Atom{P}{x}
    \land \lforall[y][(y < x \lif \lnot \Atom{P}{y})])])].
\]
यह कथन करता है कि ``समुच्चय'' की पहचान ``एकपदी संबंध'' से करने पर प्रत्येक
अरिक्त समुच्चय का एक न्यूनतम अवयव है। एक अन्य उदाहरण में, आरेखों के लिए
संयोजकता की धारणा यह कहकर व्यक्त की जा सकती है कि शीर्षों का असंयुक्त भागों
में कोई अतुच्छ विभाजन नहीं है:
\[
\lnot \lexists[A][(\lexists[x][A(x)] \land \lexists[y][\lnot A(y)]
  \land \lforall[w][\lforall[z][((\Atom{A}{w} \land \lnot \Atom{A}{z})
      \lif \lnot \Atom{R}{w,z})]])].
\]
एक और उदाहरण में, अभ्यास के रूप में ऐसा वर्ग परिभाषित करने का प्रयास कीजिए
जिसकी सदस्य !!{structure}s परिमित हों और जिनके !!{domain} का आकार सम हो। इससे भी अधिक
प्रभावशाली बात यह है कि वास्तविक संख्याओं का श्रेणीपरक वर्णन परिमेय संख्याओं
को समाहित करने वाले पूर्ण क्रमित क्षेत्र के रूप में दिया जा सकता है।

संक्षेप में, द्वितीय-कोटि तर्कशास्त्र प्रथम-कोटि तर्कशास्त्र से बहुत अधिक
अभिव्यंजक है। यह अच्छी खबर है; अब बुरी खबर। हम पहले ही कह चुके हैं कि पूर्ण
द्वितीय-कोटि अर्थविज्ञान के लिए पूर्ण कोई प्रभावी !!{derivation}-तंत्र नहीं
है। अच्छा हो या बुरा, प्रथम-कोटि तर्कशास्त्र के अनेक गुणधर्म यहाँ नहीं हैं,
जिनमें संहतता और लोवेनहाइम (L\"owenheim)--स्कोलेम प्रमेय सम्मिलित हैं।

दूसरी ओर यदि पूर्ण द्वितीय-कोटि अर्थविज्ञान छोड़कर दुर्बलतर अर्थविज्ञान
अपनाना स्वीकार्य हो, तो न्यूनतम द्वितीय-कोटि !!{derivation}-तंत्र इस
अर्थविज्ञान के लिए पूर्ण है। दूसरे शब्दों में, यदि $\Proves$ को ``न्यूनतम
तंत्र में सिद्ध करता है'' और $\Entails$ को ``दुर्बलतर अर्थविज्ञान में
तार्किक रूप से निष्पन्न करता है'' पढ़ें, तो दिखा सकते हैं कि $\Gamma \Entails !A$
होने पर सदा $\Gamma \Proves !A$। यदि !!{derivation}-तंत्र में
कुछ निश्चित परिग्रहण-अभिगृहीत सम्मिलित करने हों, तो अर्थविज्ञान को उन
द्वितीय-कोटि संरचनाओं तक सीमित करना होगा जो इन अभिगृहीतों को संतुष्ट करती
हैं: जैसे, यदि $\Delta$ परिग्रहण-अभिगृहीतों (संभवतः उन सभी) का समुच्चय है,
तो $\Gamma
\cup \Delta \Entails !A$ से $\Gamma \cup \Delta \Proves !A$
मिलता है। विशेषतः यदि विचाराधीन परिग्रहण-अभिगृहीतों के प्रयोग से $!A$
प्रमाण्य नहीं है, तो $\lnot !A$ का कोई ऐसा प्रतिरूप है जिसमें ये
परिग्रहण-अभिगृहीत फिर भी सत्य हैं।

यह देखने की सबसे सरल रीति कि पूर्णता प्रमेय दुर्बलतर अर्थविज्ञान के लिए
सत्य है, द्वितीय-कोटि तर्कशास्त्र को निम्न प्रकार बहु-वर्गीय तर्कशास्त्र
मानना है। एक वर्ग की व्याख्या सामान्य ``प्रथम-कोटि'' !!{domain} के रूप में
होती है, और फिर प्रत्येक $k$ के लिए हमारे पास ``पदिता $k$ वाले संबंधों'' का
एक !!a{domain} होता है। भाषा में अंतर्निर्मित संबंध-चिह्न
``$\Atom{\Obj{true}_k}{R,x_1,\dots,x_k}$'' लेते हैं, जिसका आशय यह कहना है कि
$R$, $x_1$, \dots,~$x_k$ के लिए सत्य है; यहाँ $R$, ``$k$-पदी संबंध'' वर्ग
का चर और $x_1$, \dots,~$x_k$ प्रथम-कोटि वर्ग की वस्तुएँ हैं।

इस अभिज्ञान के साथ दुर्बल द्वितीय-कोटि अर्थविज्ञान मूलतः बहु-वर्गीय
तर्कशास्त्र का सामान्य अर्थविज्ञान है; और हम पहले ही देख चुके हैं कि
बहु-वर्गीय तर्कशास्त्र को प्रथम-कोटि तर्कशास्त्र में अंतःस्थापित किया जा
सकता है। दोनों दिशाओं के अनुवादों को मान लें, तो द्वितीय-कोटि तर्कशास्त्र की
दुर्बलतर संकल्पना वास्तव में प्रथम-कोटि तर्कशास्त्र का छद्मवेशी रूप है,
जिसमें !!{domain} के भीतर उपयुक्त अभिगृहीतों द्वारा नियंत्रित ``वस्तुएँ'' और
``संबंध'' दोनों हैं।

\end{document}

